\documentclass[11pt]{article}
\usepackage[T1]{fontenc}
\usepackage[utf8]{inputenc}
\usepackage[ngerman]{babel}
\usepackage{amsmath,amssymb,mathtools,array}
\usepackage{geometry}
\usepackage{microtype}
\geometry{a4paper,margin=3cm}
\setlength{\parindent}{1.25em}
\setlength{\parskip}{0pt}
\newcommand{\widebar}[1]{\overline{#1}}
\newcommand{\A}{\mathcal A}
\allowdisplaybreaks
\begin{document}

\begin{center}
{\Large\bfseries 1. Über die Bildung des Formensystems der ternären\\ biquadratischen Form}\par
\vspace{1.5em}
Sitz. Ber. d. Physikal.-mediz. Sozietät in Erlangen 39 (1907), S. 176--179
\end{center}

\vspace{4em}
\begin{center}
(Auszug aus der Dissertation der Verfasserin.)
\end{center}

Mit dem Formensystem der ternären biquadratischen Form beschäftigen sich Arbeiten von Gordan, Maisano und Pascal\footnote{P. Gordan, Über das volle Formensystem der ternären biquadratischen Form $f=x_1^3x_2+x_2^3x_3+x_3^3x_1$ (Math. Annalen Bd. XVII, S. 217 bis 233. 1880); G. Maisano, 1. Sistemi completi dei primi cinque gradi della forma ternaria biquadratica e degl' invarianti, covarianti e contravarianti di sesto grado (Giorn. di Battaglini XIX). 2. Sui covarianti indipendenti di $6^\circ$ grado nei coefficienti della forma biquadratica ternaria (Rend. Circ. Mat. di Palermo I. 1887); E. Pascal, Contributo alla teoria della forma ternaria biquadratica e delle sue varie decomposizioni in fattori (Memoria premiata dalla R. Accademia delle scienze fisiche e matematiche di Napoli. 1905).}. Herr Gordan stellt das vollständige, aus 54 Bildungen bestehende, Formensystem der speziellen automorphen Form: $f=x_1^3x_2+x_2^3x_3+x_3^3x_1$ unter Zugrundelegung ähnlicher Prinzipien auf, wie er sie für die Formensysteme im binären Gebiet gegeben hat.

Bei Herrn Maisano sind für die allgemeine biquadratische Form die Formen bis zur 5. Ordnung\footnote{Unter „Ordnung“ soll die Dimension in den Koeffizienten, unter „Grad“ die in den Variablen verstanden werden.} einschließlich aufgestellt, sowie einige Invarianten, Kovarianten und Kontravarianten höherer Ordnung, nach der von Herrn Gordan in Band I der Math. Annalen für die ternäre kubische Form angewandten Methode. Herr Pascal beschäftigt sich, unter Benützung der Maisanoschen Resultate hauptsächlich mit der Frage nach dem Zerfallen der biquadratischen Form in Faktoren\footnote{In der zweiten kurzen Note versucht Herr Maisano eine lineare Abhängigkeit der 3 Kovarianten 6. Ordnung und 6. Grades nachzuweisen. Im Gegensatz zu diesem nicht ganz vollständigen Beweis glaubt Herr Pascal die lineare Unabhängigkeit der 3 Kovarianten 6. Ordnung, ebenso wie diejenige der 3 Invarianten 9. Ordnung bewiesen zu haben. Daß aber in der Tat eine lineare Relation zwischen den 3 Kovarianten einerseits, den 3 Invarianten andererseits existiert, hat sich im Laufe meiner Rechnungen ergeben; und zwar lauten die expliziten Formeln in der Bezeichnungsweise Pascal's:
\[
0=20\Omega_1+6\Omega_2-3\Omega_3+4fC_1-12fC_2-4A\cdot\Delta.
\]
\[
0=4A^3-15AB+30C-90D+9E.
\]}. 

\textbf{Das Ziel meiner Untersuchungen ist die Aufstellung des Formensystems für die allgemeine ternäre biquadratische Form; und zwar werden zunächst nur die Hauptgrundlagen gegeben, ein sogenanntes „relativ vollständiges System“\footnote{Für die Bezeichnung vgl. Gordan-Kerschensteiner, Vorlesungen über Invariantentheorie, Bd. II, S. 227.} aufgestellt.}

Der Grundgedanke der Systembildung ternärer Formen ist derselbe wie im binären Gebiet. Ausgehend von einem ersten relativ vollständigen System --- dem System der aus dem Binären übernommenen Formen --- gelangt man nach bestimmter Gesetzmäßigkeit zu Systemen mit immer höherem Modul, solange bis das System eines Moduls --- der Modul als Grundform genommen --- endlich und bekannt wird, oder auch bis ein Modul sich reduzieren läßt auf Formen, die Invarianten zum Faktor haben. Durch Überschiebung des relativ vollständigen Systems über das System des Moduls entsteht im ersten Fall das absolut vollständige System, während im zweiten Fall relativ vollständiges und absolut vollständiges System identisch werden. Infolge der durch Herrn Hilbert allgemein bewiesenen Endlichkeit der Formensysteme muß dies Verfahren notwendig zu einem Abschluß führen.

In unserm Fall läßt sich der Modul $(abc)$ des ersten relativ vollständigen Systems zurückführen auf die Moduln
\[
\Delta=(abc)^2a_x^2b_x^2c_x^2\quad\text{und}\quad \nu=(abu)^4;
\]
daraus ergibt sich leicht das relativ vollständige System mod $\nu$. Als Reihe der Moduln wählen wir nun die Formen:
\[
\nu;\qquad \nu^{(\nu)}=(\nu\nu_1x)^4=s_x^4,
\]
\[
\nu^{(s)}=(ss'u)^4\ldots
\]
und adjungieren als weitere Moduln zwei bei der Bildung des relativ vollständigen Systems mod $s$ auftretende quadratische Formen, $u_\sigma^2$ und $t_x^2$. Es läßt sich dann zeigen, daß der auf $s$ folgende Modul $(ss'u)^4$ reduzibel ist auf den Modul $(\varrho,t)$. Da aber das simultane System zweier quadratischer Formen endlich und bekannt ist, ist damit der oben gekennzeichnete Abschluß erreicht. Als relativ vollständiges System mod $(\varrho,t)$ ergeben sich 331 Bildungen. ---

Noch mache ich einige Angaben bezüglich der angewandten Methode.

Der Grundprozeß zur Erzeugung von Formen, der Faltungsprozeß, läßt sich im ternären Gebiet folgendermaßen definieren:

Ersetzt man in einem symbolischen Produkt
\[
s_x^m t_x^n u_\sigma^\mu u_\tau^\nu
\]
die Faktorenpaare
\[
\begin{array}{c|c|c|c|c}
 & s_xt_x & u_\sigma u_\tau & s_xu_\sigma\ \text{oder}\ s_xu_\tau & t_xu_\tau\ \text{oder}\ t_xu_\sigma\\
\text{resp. durch} & (stu) & (\sigma\tau x) & s_\sigma\ \text{oder}\ s_\tau & t_\tau\ \text{oder}\ t_\sigma\\
\text{Faltung} & \mathrm{I} & \mathrm{II} & \mathrm{III} & \mathrm{IV}
\end{array}
\]
so sind die entstehenden Formen durch Faltung aus der ursprünglichen hervorgegangen\footnote{Gordan. Math. Annalen, Bd. XVII, S. 219.}.

Dabei kann noch stets: $s_\sigma=0;\ t_\tau=0$ gesetzt werden. Für den Zusammenhang der einzelnen Faltungen gilt dabei der Satz:

\textbf{Die Faltungen I und II sind Grundfaltungen, aus denen sich, unabhängig von der Reihenfolge der Zusammensetzung, die Faltungen III und IV zusammensetzen lassen. In anderen Worten: Um alle aus einem gegebenen Ausdruck durch Faltung entstehenden Formen zu bilden, hat man nur die Faltungen I und II anzuwenden}\footnote{Eine Ausnahme erleidet der Satz in dem Fall, wo $n$ und $\mu$, resp. $m$ und $\nu$ gleichzeitig verschwinden; die „Formenreihe“ reduziert sich dann auf das Diagonalglied.}.

Als „Formenreihe“ definieren wir eine Anfangsform mit allen durch Faltung in sich aus dieser hervorgehenden Formen. Nach dem Satze läßt sich eine Formenreihe $s_x^m t_x^n u_\sigma^\mu u_\tau^\nu$ in ein rechteckiges Schema anordnen. Man erhält beim Fortschreiten:
\begin{enumerate}
\item um eine Kolonne nach rechts, alle durch Faltung I aus den nebenstehenden Formen entstehenden Formen,
\item um eine Zeile nach unten, alle durch Faltung II aus den obenstehenden Formen entstehenden Formen, und somit alle durch Faltung entstehenden Formen.
\end{enumerate}

Unter diesen Voraussetzungen lassen sich die Sätze über Reduzenten in ihrer allgemeinsten Form und unter der Definition
\begin{center}
„Ein Reduzent ist eine reduzible Formenreihe“
\end{center}
so aussprechen:

\begin{center}
\textbf{Ist die Anfangsform einer Formenreihe reduzibel dadurch, daß eines ihrer Glieder durch Faltung mit einem Reduzenten hervorgegangen ist, und ist die Schlußform der Formenreihe aus eben diesem Glied durch Faltung entstanden, so ist die Gesamtformenreihe reduzibel.}
\end{center}

Als weitere Reduktionsmethoden sind zu nennen:
\begin{enumerate}
\item Die sogenannte „doppelte Reduktion“; d. h. die Reduktion einer Form auf doppelte Art zur Erzielung einer Relation zwischen den höheren Formen,
\item die „Faltung mit zerfallenden Formen“, die teilweise den Charakter der Reduktion durch Reduzenten, teilweise den der „doppelten Reduktion“ zeigt.
\end{enumerate}
\clearpage

\begin{center}
{\Large\bfseries 2. Über die Bildung des Formensystems der ternären\\ biquadratischen Form}\par
\vspace{1.5em}
Journal f. d. reine u. angew. Math. 134 (1908), S. 23--90 u. zwei Tabellen
\end{center}

\vspace{3em}

\begin{center}
\textbf{Einleitung.}
\end{center}

Mit dem Formensystem der ternären biquadratischen Form beschäftigen sich Arbeiten von Gordan, Maisano und Pascal\footnote{Ein kurzer Auszug aus Einleitung und Kapitel I ist in den „Sitzungsber. der phys.-med. Sozietät Erlangen 1907“, S. 176 bis 179, erschienen.}\footnote{P. Gordan, über das volle Formensystem der ternären biquadratischen Form
\[
f=x_1^3x_2+x_2^3x_3+x_3^3x_1.
\]
(Math. Annalen, Bd. XVII (1880), S. 217--233.)
G. Maisano, 1) \emph{Sistemi completi dei primi cinque gradi della forma ternaria biquadratica e degl' invarianti, covarianti e contravarianti di sesto grado}. (Giorn. di Battaglini XIX.)
2) \emph{Sui covarianti indipendenti di $6^\circ$ grado nei coefficienti della forma biquadratica ternaria}. (Rend. Circ. Mat. di Palermo I, 1887.)
E. Pascal, \emph{Contributo alla teoria della forma ternaria biquadratica e delle sue varie decomposizioni in fattori}. (Memoria premiata dalla R. Accademia delle scienze fisiche e matematiche di Napoli. 1905.)}. Herr Gordan stellt das vollständige, aus 54 Bildungen bestehende, Formensystem der speziellen automorphen Form:
\[
f=x_1^3x_2+x_2^3x_3+x_3^3x_1
\]
unter Zugrundelegung ähnlicher Prinzipien auf, wie er sie für die Formensysteme im binären Gebiet gegeben hat.

Bei Herrn Maisano sind für die allgemeine biquadratische Form die Formen bis zur 5. Ordnung\footnote{Unter „Ordnung“ soll die Dimension in den Koeffizienten, unter „Grad“ die in den Variablen verstanden werden.} einschließlich aufgestellt, sowie einige Invarianten, Kovarianten und Kontravarianten höherer Ordnung, nach der von Herrn Gordan in Band I der Math. Annalen für die ternäre kubische Form angewandten Methode. Herr Pascal beschäftigt sich, unter Benutzung der Maisanoschen Resultate, hauptsächlich mit der Frage nach dem Zerfallen der biquadratischen Form in Faktoren\footnote{In der zweiten kurzen Note versucht Herr Maisano eine lineare Abhängigkeit der drei Kovarianten 6. Ordnung und 6. Grades nachzuweisen. Im Gegensatz zu diesem nicht ganz vollständigen Beweis glaubt Herr Pascal die lineare Unabhängigkeit der drei Kovarianten 6. Ordnung, ebenso wie diejenige der drei Invarianten 9. Ordnung bewiesen zu haben. Daß aber in der Tat eine lineare Relation zwischen den drei Kovarianten einerseits, den drei Invarianten andererseits existiert, zeigen die Formeln (13), § 11 und (22), § 17 Anmerkung, der vorliegenden Arbeit, wo diese Relationen explizit gegeben sind. Nach einer Mitteilung des Herrn Pascal erklärt sich der Widerspruch durch einen numerischen Fehler im Ausdruck seiner Kontravariante $p$ (hier $\varrho$ genannt) S. 46 seiner Abhandlung.}.

\emph{Das Ziel der folgenden Untersuchungen ist die Aufstellung des Formensystems für die allgemeine ternäre biquadratische Form; und zwar werden in dieser Arbeit nur die Hauptgrundlagen gegeben, ein sogenanntes „relativ vollständiges System“\footnote{Für die Bezeichnung vgl. § 3a.} aufgestellt.}

Die Arbeit schließt sich eng an die Gordansche Arbeit an; doch waren die dort nur ganz im allgemeinen gegebenen Prinzipien erst im einzelnen auszuarbeiten. Mit Hilfe eines Satzes über den Zusammenhang der Faltungen und durch Einführung der „Formenreihe“ (§ 1) werden die Reduktionssätze scharf formuliert und vervollständigt (§ 3), während die rekurrierende Aufstellung spezieller Reihenentwicklungen (§ 2) das rechnerische Mittel zur wirklichen Durchführung der Reduktionen gibt.

Der Grundgedanke der Systembildung ternärer Formen ist derselbe wie im binären Gebiet. Ausgehend von einem ersten relativ vollständigen System --- dem System der aus dem Binären übernommenen Formen ---, gelangt man nach bestimmter Gesetzmäßigkeit zu Systemen mit immer höherem Modul, solange bis das System eines Moduls --- der Modul als Grundform genommen --- endlich und bekannt wird, oder auch bis ein Modul sich reduzieren läßt auf Formen, die Invarianten zum Faktor haben. Durch Überschiebung des relativ vollständigen Systems über das System des Moduls entsteht im ersten Fall das absolut vollständige System, während im zweiten Fall relativ vollständiges und absolut vollständiges System identisch werden. Infolge der durch Herrn Hilbert allgemein bewiesenen Endlichkeit der Formensysteme muß dies Verfahren notwendig zu einem Abschluß führen.

In unserm Fall wird der Modul $(abc)$ des ersten relativ vollständigen Systems (§ 4) zurückgeführt auf die Moduln
\[
\Delta=(abc)^2a_x^2b_x^2c_x^2\quad\text{und}\quad \nu=(abu)^4
\]
(§ 5), und sodann das relativ vollständige System mod $\nu$ gefunden (§ 6). Diese Resultate sind schon in der Gordanschen Arbeit für die allgemeine biquadratische Form gegeben; unsere Art der Ableitung läßt sich jedoch leichter auf Formen höheren Grades übertragen.

Als Reihe der Moduln wählen wir nun die Formen (§ 7):
\[
\nu,\qquad \nu(\nu)=(\nu\nu_1x)^4=s_x^4,\qquad \nu(s)=(ss'u)^4,\ldots .
\]

Bei der Bildung des relativ vollständigen Systems mod $s$ werden zwei im System auftretende quadratische Formen, $u_\varrho^2$ und $t_x^2$, als Moduln adjungiert (§ 9 und 10) und demzufolge das relativ vollständige System mod $(s,\varrho,t)$ gebildet; es wird sodann gezeigt (§ 17), daß der Modul $(ss'u)^4$ reduzibel ist auf die Moduln $\varrho$ und $t$. Dadurch geht das nächsthöhere System über in ein relativ vollständiges System mod $(\varrho,t)$. Da aber das simultane System zweier quadratischen Formen endlich und bekannt ist, im allgemeinen Fall aus 20 Bildungen besteht\footnote{Vgl. Clebsch-Lindemann, Vorlesungen über Geometrie. (Bd. I, S. 288 ff.) System zweier kogredienter Formen. Gordan, Über Büschel von Kegelschnitten. (Math. Ann. Bd. XIX. S. 530). System zweier kontragredienter Formen.}, ist mit der Aufstellung des relativ vollständigen Systems mod $(\varrho,t)$ der oben gekennzeichnete Abschluß erreicht. Die Überschiebung dieses Systems über das System von $(\varrho,t)$ zur Bildung des absolut vollständigen Systems bleibt vorbehalten.

Als relativ vollständiges System mod $(\varrho,t)$ werden 331 Bildungen gefunden, die in der beigefügten Tabelle nach ihrem Grad in den Variablen $x$ und $u$ geordnet sind.

Zum Schluß geben wir noch eine Übersicht über die eingeführten Symbole:
\begin{align*}
f&=a_x^4,\\
\theta&=\theta_x^4u_\vartheta^2=(abu)^2a_x^2b_x^2,\\
K&=K_x^6u_k^3=(a\theta u)u_\vartheta^2a_x^3\theta_x^3,\\
j&=u_j^6=(a\theta u)^4u_\vartheta^2,\\
\Delta&=\Delta_x^6=a_\vartheta^2a_x^2\theta_x^4=(abc)^2a_x^2b_x^2c_x^2,\\
N&=N_x^8u_\eta=(a\Delta u)a_x^3\Delta_x^5,\\
\nu&=u_\nu^4=(abu)^4,\\
H&=H_x^2u_\eta^4=(\nu\nu_1x)^2u_\nu^2u_{\nu_1}^2,\\
L&=L_x^3u_l^6=(\nu\eta x)H_x^2u_\nu^3u_\eta^3,\\
g&=g_x^6=(\nu\eta x)^4H_x^2,\\
\sigma&=u_\sigma^6=H_\nu^2u_\nu^2u_\eta^4,\\
s&=s_x^4=(\nu\nu_1x)^4,\\
Z&=Z_x^4u_\zeta^2=(ss'u)^2s_x^2s_x'^2,\\
\varrho&=u_\varrho^2=\theta_\nu^4u_\vartheta^2,\\
t&=t_x^2=a_\eta^4H_x^2,\\
i&=a_\nu^4,\\
J&=s_\nu^4.
\end{align*}

\begin{center}
{\Large\bfseries Kapitel I. \quad Allgemeine Sätze über ternäre Formen.}
\end{center}

\begin{center}
\textbf{§ 1. \quad Faltungsprozeß. \quad Formenreihen.}
\end{center}

Der Grundprozeß zur Erzeugung von Formen ist der Faltungsprozeß, der sich im ternären Gebiet folgendermaßen definieren läßt. Gegeben sei ein symbolisches Produkt:
\[
s_x^m t_x^n u_\sigma^\mu u_\tau^\nu .
\]
Ersetzt man die Faktorenpaare
\[
s_xt_x\qquad u_\sigma u_\tau\qquad s_xu_\sigma\ \text{oder}\ s_xu_\tau
\qquad t_xu_\tau\ \text{oder}\ t_xu_\sigma
\]
bzw. durch
\[
(stu)\qquad (\sigma\tau x)\qquad s_\sigma\ \text{oder}\ s_\tau
\qquad t_\tau\ \text{oder}\ t_\sigma,
\]
Faltung
\[
\text{I}\qquad\text{II}\qquad\text{III}\qquad\text{IV},
\]
so sind die entstehenden Formen durch Faltung aus der ursprünglichen hervorgegangen\footnote{Gordan, Math. Annalen Bd. XVII S. 219.}.

Dabei läßt sich der Ausdruck $s_x^m t_x^n u_\sigma^\mu u_\tau^\nu$ entweder als wirkliches Produkt zweier Formen $S\cdot T$ auffassen, oder auch als einzige Form mit mehreren Reihen von Symbolen:
\[
s_x^m t_x^n u_\sigma^\mu u_\tau^\nu=A_x^{m+n}u_x^{\mu+\nu}.
\]
Im ersten Fall sprechen wir von der Faltung der Form $S$ mit $T$, im zweiten Fall von der „Faltung der Form $A$ in sich“. Der Kürze halber nehmen wir im folgenden immer an (was in der Tat bei allen später zu betrachtenden Formen der Fall ist)
\[
\tag{a}
s_\sigma^\lambda=0,\qquad t_\tau^\lambda=0
\]
für beliebiges von 0 verschiedenes $\lambda$ und $x$ und vereinigt mit beliebigen anderen Faltungen. Es sagt dies aus, daß die Form $s_x^m t_y^n u_\sigma^\mu u_\tau^\nu$ eine sogenannte „Normalform“ ist, d. h. bei Anwendung des $\Omega$-Prozesses, sowohl nach den Variablen $x$ und $u$, wie nach den Variablen $y$ und $v$, verschwindet.

Die Bedingung (a) stellt also keine Einschränkung dar, sondern läßt sich stets erreichen\footnote{Vgl. Gordan, Math. Annalen Bd. V S. 104.}.

Für den Zusammenhang der einzelnen Faltungen gilt nun folgender Satz:

\emph{Satz I. Die Faltungen I und II sind Grundfaltungen, aus denen sich, unabhängig von der Reihenfolge der Zusammensetzung, die Faltungen III und IV zusammensetzen lassen. In anderen Worten: Um alle aus einem gegebenen Ausdruck durch Faltung entstehenden Formen zu bilden, hat man nur die Faltung I und II anzuwenden.}

Beweis: Nach dem Identitätssatz, bzw. Produktsatz für Matrizen, folgt unter Berücksichtigung von (a):
\[
(\widehat{st}\sigma x)=-t_\sigma,\qquad
(\widehat{\sigma\tau}su)=-s_\tau,
\]
\[
(\widehat{st}\tau x)=s_\tau,\qquad
(\widehat{\sigma\tau}tu)=t_\sigma,
\]
\[
(stu)(\sigma\tau x)=s_\tau+t_\sigma-u_x\,s_\tau t_\sigma .
\]
Durch Zusammensetzung von Faltung I und II entstehen also Faltung III und IV, und zwar ebenso bei Anwendung von Faltung II auf Faltung I, d. h. auf die Faktoren $(stu)u_\sigma u_\tau$, wie bei Anwendung von Faltung I auf Faltung II, auf die Faktoren $(\sigma\tau x)s_xt_x$.

Als \emph{Formenreihe}\footnote{Vgl. Clebsch, Über eine Fundamentalaufgabe der Invariantentheorie (Abh. der Gött. Ges. d. Wiss. Bd. XVII): § 17 und § 18 Schluß. Die „Formenreihe“ unterscheidet sich von dem dort eingeführten „eigentlich reduzierten äquivalenten System“ durch das Prinzip der Anordnung nach höheren Formen.} definieren wir eine Anfangsform mit allen durch Faltung in sich aus dieser hervorgegangenen Formen und bezeichnen die Formenreihe durch die Anfangsform. Als „höhere Form“ soll hierbei jede Form mit mehr Faltung in sich gelten, während unter den gleichberechtigten Faltungen $s_\tau$ und $t_\sigma$ eine als höher normiert werden muß. Nach dem Bildungsgesetz der Formenreihe folgt:

1) Jede Form der Formenreihe ist linear in den Symbolen der Anfangsform.

2) Die Formenreihe ist bei Anfangsformen mit je zwei Reihen kontragredienter Symbole eindeutig bestimmt; bei Anfangsformen mit mehr Symbolreihen ist sie eindeutig zu normieren durch Auszeichnung spezieller unter den durch den Identitätssatz verbundenen Faltungen.

Nach Satz I. läßt sich eine Formenreihe $s_x^m t_x^n u_\sigma^\mu u_\tau^\nu$ ($m>n;\ \mu>\nu$) nach folgendem rechteckigen Schema anordnen:
\begingroup\small
\[
\begin{array}{c|c|c|c}
s_x^m t_x^n u_\sigma^\mu u_\tau^\nu
& (stu)
& (stu)^2
& \cdots\\[0.3em]
(\sigma\tau x)
& t_\sigma,\ s_\tau
& t_\sigma(stu),\ s_\tau(stu)
& \cdots\\[0.3em]
(\sigma\tau x)^2
& t_\sigma(\sigma\tau x),\ s_\tau(\sigma\tau x)
& t_\sigma^2,\ t_\sigma s_\tau,\ s_\tau^2
& \cdots\\[0.3em]
\vdots&\vdots&\vdots&\ddots\\[0.3em]
(\sigma\tau x)^\nu
& t_\sigma(\sigma\tau x)^{\nu-1},\ s_\tau(\sigma\tau x)^{\nu-1}
& t_\sigma^2(\sigma\tau x)^{\nu-2},\ t_\sigma s_\tau(\sigma\tau x)^{\nu-2},\ s_\tau^2(\sigma\tau x)^{\nu-2}
& \cdots\\[0.3em]
\vdots
& t_\sigma(\sigma\tau x)^\nu
& t_\sigma^2(\sigma\tau x)^{\nu-1},\ t_\sigma s_\tau(\sigma\tau x)^{\nu-1}
& \cdots
\end{array}
\]
\endgroup
und der entsprechend fortgesetzte untenstehende Teil\footnote{Hier, und ähnlich im folgenden, ist der mit einem Sternchen bezeichnete unten stehende Teil an die ebenso bezeichnete rechts oben befindliche Stelle des Schemas anzusetzen.}
\begingroup\small
\[
\begin{array}{c|c|c}
(stu)^n & \cdots & \cdots\\[0.3em]
t_\sigma(stu)^{n-1},\ s_\tau(stu)^{n-1}
& s_\tau(stu)^n & \cdots\\[0.3em]
t_\sigma^2(stu)^{n-2},\ t_\sigma s_\tau(stu)^{n-2},\ s_\tau^2(stu)^{n-2}
& t_\sigma s_\tau(stu)^{n-1},\ s_\tau^2(stu)^{n-1}
& s_\tau^2(stu)^n .
\end{array}
\]
\endgroup

Man erhält hier, wenn man

1) um eine Kolonne nach rechts fortschreitet, alle durch Faltung I aus den nebenstehenden Formen entstehenden Formen,

2) um eine Zeile nach unten fortschreitet, alle durch Faltung II aus den obenstehenden Formen entstehenden Formen,

und somit alle durch Faltung entstehenden Formen.

Eine beliebige Form der Gesamtformenreihe: $t_\sigma^\kappa s_\tau^\lambda(stu)^\mu$ oder $t_\sigma^\kappa s_\tau^\lambda(\sigma\tau x)^\nu$ bildet den Anfang einer neuen Formenreihe, die aus all denjenigen Formen des rechteckigen Schemas mit der Form $t_\sigma^\kappa s_\tau^\lambda(stu)^\mu$ oder $t_\sigma^\kappa s_\tau^\lambda(\sigma\tau x)^\nu$ als linkem oberem Eckpunkt besteht, die den Faktor $t_\sigma^\kappa s_\tau^\lambda$ haben.

Nachbemerkung. Eine Ausnahme erleidet Satz I in dem Fall, wo die Zahlen $n$ und $\mu$ (bzw. $m$ und $\nu$) gleich Null werden, Faltung I, II und IV also nicht mehr existieren, wohl aber Faltung III (bzw. IV). Als Formenreihe bezeichnen wir in diesem Fall das Diagonalglied des Schemas:
\[
\begin{array}{ccccc}
s_x^m u_\tau^\nu &&&& t_x^n u_\sigma^\mu\\
& s_\tau && t_\sigma &\\
& s_\tau^2 && t_\sigma^2 &\\
&&\ddots&&\\
&&s_\tau^\nu && t_\sigma^\mu .
\end{array}
\]
In Übereinstimmung mit dem Fehlen von Faltung I und II reduziert sich in diesem Fall die Reihenentwicklung nach Polaren der Formenreihe stets auf das Anfangsglied; es gelten aber die Sätze über Reduzenten.

\begin{center}
\textbf{§ 2. \quad Reihenentwicklungen nach Polaren der Formenreihe.}
\end{center}

Für Formen mit $n$ Variablen sind zwei Arten von Reihenentwicklungen explizit aufgestellt\footnote{Vgl. Gordan, Über Kombinanten. § 2 und 5 (Math. Ann. Bd. V).}:

1) für Formen mit je einer Reihe kontragredienter Variablen: $s_x^m u_\tau^\nu$ eine nach Potenzen von $u_x$ fortschreitende Reihe, deren Koeffizienten „Normalformen“ sind;

2) für Formen mit zwei Reihen kogredienter Variablen $s_x^m t_x^n$ eine Entwicklung nach Polaren der Elementarkovarianten (Polaren der Formenreihe $s_x^m t_x^n$), die ternär folgendermaßen lautet\footnote{Vgl. Study, Methoden zur Theorie der ternären Formen (Teubner 1889) II. § 7. Unsere Ableitung gibt im Unterschied von der dortigen die Methode zur raschen rechnerischen Bestimmung der numerischen Koeffizienten $C_{pr\varrho}$. Die Clebsch-Studysche Entwicklung würde bei Spezialisierung a), a') lauten:
\[
(stu)^\mu u_\tau^{\nu-\kappa}v_\tau^\kappa s_x^m t_x^{\,n-\lambda}(tuv)^\lambda
=\sum_{p,r,\varrho} C_{pr\varrho}
\bigl[s_\tau^\varrho(stu)^{\mu+r-\varrho}\bigr]_{\widehat{uv}^{\lambda-r+\varrho}v^{\nu-\varrho},\,u^p,\,(vw=\widehat{xw})},
\]
läßt also nicht die Vereinfachung zu, die bei unserer Entwicklung schon im Laufe der Rechnung eintritt, wenn man weiß, daß alle Glieder mit dem Faktor $u_x$ auszulassen sind.}:
\[
\tag{I}
s_x^m t_x^{n-\lambda}t_y^\lambda
=\sum_{\varrho}
\frac{\binom{m}{\varrho}\binom{\lambda}{\varrho}}
{\binom{m+n-\varrho+1}{\varrho}}
\bigl[(st\widehat{xy})^\varrho s_x^{m-\varrho}t_x^{n-\varrho}\bigr]_{y^{\lambda-\varrho}}.
\]

Wir kombinieren beide Entwicklungen, indem wir für Formen mit je zwei Reihen kontragredienter Variablen:
\[
s_x^m t_x^{n-\lambda}t_y^\lambda u_\sigma^\mu u_\tau^{\nu-\kappa}v_\tau^\kappa
\]
Entwicklungen nach Potenzen von $u_x$ aufstellen, deren Koeffizienten zusammengesetzte Polaren der Formenreihe $s_x^m t_x^n u_\sigma^\mu u_\tau^\nu$ sind, genommen nach den Variablen $x$ und $u$.

Es genügt, die Reihe für je zwei kontragrediente Symbol- (bzw. Variablen-) reihen aufzustellen, da sich durch Zusammenfassen von je zwei Symbol- (Variablen-) reihen zu einer einzigen neuen Reihe Formen mit mehr Symbol- (Variablen-) reihen auf diesen Fall zurückführen lassen.

Um zu erreichen, daß alle Formen der Formenreihe nur je zwei Symbol- bzw. Variablenreihen enthalten, spezialisieren wir:
\[
\begin{array}{ll}
\text{a) } \sigma=\widehat{st}, & \text{a') } y=\widehat{uv},\\[0.3em]
\text{b) } s=\widehat{\sigma\tau}, & \text{b') } v=\widehat{xy},
\end{array}
\]
und führen die Entwicklung für die Spezialisierung a), a') durch.

Durch direkte Übertragung der Entwicklung I erhält man zusammengesetzte Polaren für Formen $(stu)^\mu s_x^m t_x^{n-\lambda}t_y^\lambda$, während für Formen $(stu)^\mu u_\tau^{\nu-\kappa}v_\tau^\kappa s_x^m t_x^{n-\lambda}t_y^\lambda$ anstatt zusammengesetzter Polaren Glieder solcher entstehen, deren Auswertung die Einführung immer neuer, erst am Schlusse gleichzusetzender Hilfsvariablen nötig macht, wodurch die Rechnung unnötig kompliziert wird.

Wir schlagen deshalb einen indirekten Weg ein, indem wir die zusammengesetzten Polaren aller Formen der Formenreihe, also Ausdrücke
\[
\bigl[s_\tau^\varrho(stu)^{\mu-\varrho+r}\bigr]_{y^\lambda}^{v^\kappa}
\]
darstellen durch die Anfangsglieder dieser Polaren und durch Umkehrung des entstehenden rekurrierenden Gleichungssystems die gewünschte Entwicklung nach Polaren erhalten.

Nach dem Übertragungsprinzip gilt für die $\lambda$-te Polare einer Form $s_x^m t_x^n:[s_x^m t_x^n]_{y^\lambda}$ ($\lambda\le n$) folgende Entwicklung (der Fall $\lambda>n$ führt durch die Entwicklung von $[s_y^m t_y^n]_{x^{m+n-\lambda}}$ auf den ersten zurück)\footnote{Gordan, Die Resultante binärer Formen (Kap. I, § 5). Rend. Circ. Mat. di Palermo XXII. 1906.}:
\[
\bigl[s_x^m t_x^n\bigr]_{y^\lambda}
=\sum_r (-1)^r\,
\frac{\binom{m}{r}\binom{\lambda}{r}}{\binom{m+n}{r}}\,
(st\widehat{xy})^r s_x^{m-r}t_x^{n-\lambda}t_y^{\lambda-r},
\]
und entsprechend für die $\kappa$-te Polare von $u_\sigma^\mu u_\tau^\nu:[u_\sigma^\mu u_\tau^\nu]_{v^\kappa}$
\[
\bigl[u_\sigma^\mu u_\tau^\nu\bigr]_{v^\kappa}
=\sum_\varrho (-1)^\varrho\,
\frac{\binom{\mu}{\varrho}\binom{\kappa}{\varrho}}{\binom{\mu+\nu}{\varrho}}\,
(\sigma\tau\widehat{uv})^\varrho u_\sigma^{\mu-\varrho}u_\tau^{\nu-\kappa}v_\tau^{\kappa-\varrho}.
\]

Durch Multiplikation folgt, unter Einführung von Spezialisierung a) und a')
\[
\begin{aligned}
\bigl[s_x^m t_x^n(stu)^\mu u_\tau^\nu\bigr]_{\widehat{uv}^{\lambda}}^{v^\kappa}
={}&
\sum_{r,\varrho}c_{r\varrho}\,
(st\widehat{uv})^r(st\widehat{uv})^\varrho
s_x^{m-r}t_x^{n-\lambda}(tuv)^{\lambda-r}(stu)^{\mu-\varrho}
u_\tau^{\nu-\kappa}v_\tau^{\kappa-\varrho},
\end{aligned}
\]
und daraus, durch Entwicklung nach Potenzen von $u_x$ unter Auszeichnung des ersten Gliedes ($\sum'_{r,\varrho}$ bedeutet, daß das Glied $r=0$, $\varrho=0$ auszulassen ist) und Berücksichtigung von (a) S. 26
\[
\tag{II}
\begin{aligned}
&s_x^m t_x^{n-\lambda}(tuv)^\lambda(stu)^\mu u_\tau^{\nu-\kappa}v_\tau^\kappa\\
&\quad =
\bigl[s_x^m t_x^n(stu)^\mu u_\tau^\nu\bigr]_{\widehat{uv}^{\lambda}}^{v^\kappa}\\
&\qquad
-\sum_{r,\varrho}' c_{r\varrho}\,s_\tau^\varrho(stu)^{\mu-\varrho+r}
u_\tau^{\nu-\kappa}v_\tau^{\kappa-\varrho}
s_x^{m-r}t_x^{n-\lambda}(tuv)^{\lambda-r+\varrho}v_x^r\\
&\qquad
+u_x\sum_{\varrho}\sum_{r=1}^{\lambda} c_{r\varrho}\,
s_\tau^\varrho(stu)^{\mu-\varrho+r-1}(stv)
u_\tau^{\nu-\kappa}v_\tau^{\kappa-\varrho}
s_x^{m-r}t_x^{n-\lambda}(tuv)^{\lambda-r+\varrho}v_x^{r-1}\\
&\qquad
\vdots\\
&\qquad
-(-1)^\lambda u_x^\lambda\sum_\varrho c_{\lambda\varrho}\,
s_\tau^\varrho(stu)^{\mu-\varrho}(stv)^\lambda
u_\tau^{\nu-\kappa}v_\tau^{\kappa-\varrho}
s_x^{m-\lambda}t_x^{n-\lambda}(tuv)^\varrho,
\end{aligned}
\]
wobei
\[
c_{r\varrho}=(-1)^{r+\varrho}\cdot
\frac{\binom{m}{r}\binom{\lambda}{r}}{\binom{m+n}{r}}\cdot
\frac{\binom{\mu}{\varrho}\binom{\kappa}{\varrho}}{\binom{\mu+\nu}{\varrho}}
\qquad
\begin{array}{l}
\text{für }\lambda\le n,\\
\kappa\le\nu.
\end{array}
\]
und entsprechend für die übrigen Zusammenstellungen der Zahlen $\lambda,n;\kappa,\nu$. Der Ausdruck $s_x^m t_x^{n-\lambda}(tuv)^\lambda(stu)^\mu u_\tau^{\nu-\kappa}v_\tau^\kappa$ ist also zurückgeführt auf eine zusammengesetzte Polare, und, unter Anwendung der Identität
\[
(stv)u_\tau=(stu)v_\tau-s_\tau(tuv),
\]
auf eine Summe von analog gebildeten Ausdrücken, die höheren Formen der Formenreihe entsprechen.

Durch Iteration gelangt man zu der Entwicklung:
\[
\tag{III}
s_x^m t_x^{n-\lambda}(tuv)^\lambda(stu)^\mu u_\tau^{\nu-\kappa}v_\tau^\kappa
=
\sum_{p,r,\varrho}u_x^p\,C_{pr\varrho}\cdot
\bigl[s_\tau^\varrho(stu)^{\mu-\varrho+r}\bigr]_{\widehat{uv}^{\lambda-r+\varrho}v^{\kappa-\varrho+p},\,v_x^{r-p}}.
\]
Die numerischen Koeffizienten $C_{pr\varrho}$ berechnen sich eindeutig und rekurrierend aus den Koeffizienten $c_{r\varrho}$, $c'_{r\varrho}$ des durch Iteration von (II.) entstehenden Gleichungssystems (vgl. das Beispiel S. 42). Analoge Entwicklungen ergeben sich bei den übrigen Spezialisierungen.

\begin{center}
\textbf{§ 3. \quad Reduktionssätze.}
\end{center}

Die bekannten Sätze über die Reduktion von Formen und Formensystemen sollen nun mit den Paragraphen 1 und 2 in Zusammenhang gebracht werden, wozu einige aus der Theorie der binären Formen übernommene Definitionen nötig sind.

a) Wir bezeichnen als \emph{relativ vollständiges System} modulo einer vorgegebenen Reihe von Formen, ein System von Formen von der Eigenschaft, daß alle durch Faltung eines beliebigen Produktes dieser Formen entstehenden Bildungen sich ausdrücken lassen als ganze rationale Funktionen der Formen des Systems, bis auf ein additives Glied von Ausdrücken, die durch Faltung der Systemformen mit dem System des Moduls entstanden sind\footnote{Gordan-Kerschensteiner, Vorlesungen über Invariantentheorie. Bd. II. S. 227.}.

b) Wir nennen eine Form dann \emph{reduzibel}, wenn sie sich ausdrücken läßt durch Formen, die Invarianten zum Faktor haben oder durch „höhere Formen“; d. h. durch höhere Formen der Gesamtformenreihe, der die Form angehört, oder durch Formen, die die Symbole des Moduls in höherer Ordnung enthalten.

Die Sätze über Reduzenten\footnote{Gordan, Math. Annalen Bd. XVII. S. 222.} lassen sich nun in ihrer allgemeinsten Form in folgender Weise aussprechen:

Definition: \emph{Ein Reduzent ist eine reduzible Formenreihe.}

\emph{Satz II. Ist die Anfangsform einer Formenreihe reduzibel dadurch, daß eines ihrer Glieder durch Faltung mit einem Reduzenten hervorgegangen ist (einen Reduzenten zum Faktor hat), und ist die Schlußform der Formenreihe aus eben diesem Glied durch Faltung entstanden, so ist die Gesamtformenreihe reduzibel}\footnote{Die Vereinfachung, die durch Satz II eintritt, läßt sich an folgendem Beispiel zeigen: Für die spezielle Form $f=x_1^3x_2+x_2^3x_3+x_3^3x_1$ ist $a_y^2a_x^2u_y^2$ Reduzent. Daraus folgt nach Satz II die Reduktion der Formenreihe
\[
\theta_\nu(\vartheta\nu x)\theta_x^3u_\vartheta u_\nu^2
=
a_\nu^2(abu)-a_\nu b_\nu(aba),
\]
da $\theta_\nu^4=a_\nu^2b_\nu^2(abu)^2$ aus dem Glied $a_\nu^2(abu)$ durch Faltung entstanden. Bei Gordan, a. a. O. S. 231, wird hingegen die Reduktion der Formen
\[
\theta_\nu(\vartheta\nu x),\quad \theta_\nu(\vartheta\nu x)^2,\quad
\theta_\nu^2,\quad \theta_\nu^2(\vartheta\nu x),\quad
\theta_\nu^2(\vartheta\nu x)^2
\]
einzeln durchgeführt.}.

Beweis: Die Formenreihe entsteht aus dem Anfangsglied durch Faltung in sich, also durch höhere Faltung mit dem Reduzenten einerseits, durch Faltung des Reduzenten in sich andererseits. In beiden Fällen lassen sich, nach § 2, alle Formen der Formenreihe darstellen als Polaren (Überschiebungen) über die Formenreihe des Reduzenten.

Der Schluß von der Anfangsform auf die Gesamtformenreihe hat nicht mehr statt bei den übrigen Reduktionsmethoden. Es sind dies

1) die sogenannte „doppelte Reduktion“.

Wir verstehen darunter die Reduktion eines Ausdrucks, der zwei voneinander unabhängige Reduzenten zum Faktor hat, auf doppelte Art, wodurch Relationen zwischen den höheren Formen entstehen, auf die reduziert wurde.

Durch systematische Anwendung der „doppelten Reduktion“ muß man einerseits alle Relationen erhalten\footnote{So wurden beispielsweise durch „doppelte Reduktion“ gefunden: die Reduktion der Form $a_\eta^4$ (§ 10), der einzigen bei Maisano überflüssig ins System aufgenommenen Form; sodann die in der Anmerkung zur Einleitung erwähnten Relationen zwischen den 3 Kovarianten und den 3 Invarianten.}, andrerseits hat man, wenn die „doppelte Reduktion“, auf verschiedene Ausdrücke angewandt, keine neuen Relationen liefert, sowohl ein Kriterium für die Unabhängigkeit der höheren Formen wie eine Kontrolle der Rechnung.

2) \emph{Faltung mit zerfallenden Formen.}

Unter „zerfallenden Formen“ sollen --- mit geringer Verallgemeinerung des gebräuchlichen Begriffs --- Formen verstanden werden, die sich ausdrücken lassen durch Produkte von Formen niedrigeren Grades und durch „höhere“ Formen. Produkte von Formen gehen bei einmaliger Faltung in Produkte über; ebenso Produkte von nichtlinearen Formen bei zweimaliger Faltung bei den Spezialisierungen a') und b') (§ 2) nach dem Produktsatz:
\[
a_yb_y a_xb_x=\frac12a_y^2b_x^2+\frac12b_y^2a_x^2-\frac12(ab\widehat{yx})^2.
\]
Erste Polaren (Überschiebungen) und gewisse zweite Polaren zerfallender Formen sind also wieder zerfallende Formen. Es folgt:

Ein durch einmalige (bzw. zweimalige) Faltung mit einer zerfallenden Form entstandenes Glied einer einmaligen (zweimaligen) Überschiebung über eine zerfallende Form wird selbst eine zerfallende Form:

a) wenn die nächsthöheren Formen in der Formenreihe der zerfallenden Form zerfallen oder reduzibel werden, oder auch wenn bei einmaliger Faltung die Spezialisierungen $y=\widehat{uv}$ oder $v=\widehat{xy}$ eintreten;

b) wenn die übrigen einmaligen Faltungen auf höhere Formen führen (vgl. § 12, B. $H_k$ und $H_k(k\eta x)$).

Ein solches Glied einer Überschiebung kann auch zerfallen:

c) wenn die übrigen einmaligen Faltungen auf niedrigere Formen führen, die unabhängig von der Faltung mit der zerfallenden Form zerfallen, d. h. sich ausdrücken lassen durch Produkte und durch die zu reduzierende Form, wobei man sich aber jeweils erst durch Rechnung überzeugen muß, daß der numerische Koeffizient des betreffenden Gliedes nicht identisch verschwindet. Es ist dies nichts weiter als „doppelte Reduktion“ der niederen Form (vgl. § 23, C. $H_q(\theta H u)(\theta s u)$).

Zerfallende Formen entstehen durch alle einmaligen Faltungen mit Funktionaldeterminanten\footnote{vgl. Gordan, a. a. O. § 3.}, außerdem durch gewisse höhere Faltungen. Es wird
\[
(st\widehat{yx})s_y=\frac12\,t\,s_y^2-\frac12\,s\,t_y^2+\frac12(st\widehat{yx})^2,
\]
\[
(st\widehat{yx})^2s_y^2=\frac13\,t\,s_y^3-\frac13\,s\,t_y^3+s_y(st\widehat{yx})^2-\frac13(st\widehat{yx})^3.
\]
Analoge Formeln gelten für die dualistischen Formen
\[
(\sigma\tau\widehat{\nu u})v_\sigma
\quad\text{und}\quad
(\sigma\tau\widehat{\nu u})v_\sigma^2.
\]

Nach den Sätzen dieses Paragraphen ergibt sich für die Aufstellung aller irreduziblen Formen, die aus der Faltung von $S$ mit $T$ entstehen, folgende Regel:

Man gehe, bei den niedrigsten Faltungen beginnend, auf beliebigem, vorher festgelegtem Wege (z. B. zeilenweise oder symmetrisch zum Diagonalglied) so weit in der Bildung der Formenreihe $ST$ voran, bis man an eine Form gelangt, die einen Reduzenten zum Faktor hat. Die durch diese Form definierte Formenreihe ist auszulassen, sobald die Schlußform die in Satz II aufgestellte Bedingung erfüllt. Nach Ausscheidung aller reduziblen Formenreihen sind durch Anwendung der doppelten Reduktion alle übrigen reduziblen, ebenso alle zerfallenden Formen zu entfernen. Stellen der Gesamtformenreihe, die durch unabhängig voneinander reduzible Formenreihen doppelt überdeckt sind, geben Anlaß zur Reduktion von höheren Formen (vgl. § 11, D).


\clearpage
\noindent\textbf{Satz III.}
\emph{Die aus dem binären Gebiet übernommenen Formen, d.\,h. diejenigen Formen, die aus den Systemformen der entsprechenden binären Form nach dem Clebschschen Übertragungsprinzip durch Ränderung entstehen, bilden ein relativ vollständiges System mod $(abc)$.}

\noindent\textbf{Beweis:}
Einer binären Relation, die aussagt, daß die Formen $Q'_1,\ldots,Q'_n$ ein vollständiges System einer binären Grundform bilden:
\[
Q'-F(Q')=0=\sum_\lambda\bigl\{(ab)c_x+(bc)a_x+(ca)b_x\bigr\}\varphi_\lambda^{*}
\]
entspricht durch Ränderung die ternäre Relation:
\[
Q-F(Q_i)=\sum_\lambda u_x\cdot(abc)\varphi_\lambda .
\]
Dies ist aber die Definitionsgleichung für ein relativ vollständiges System mod $(abc)$. Für die biquadratische Form besteht das System der übernommenen Formen aus folgenden Formen:
\[
\begin{gathered}
f=a_x^4,\qquad
\theta=\theta_x^4u_\vartheta^2=(abu)^2a_x^2b_x^2,\qquad
K=K_x^6u_k^3=(a\theta u)u_\vartheta^2a_x^3\theta_x^3,\\
j=u_j^6=(a\theta u)^4u_\vartheta^2,\qquad
\nu=u_\nu^4=(abu)^4,
\end{gathered}
\]
unter denen die vier ersten ein relativ vollständiges System mod $((abc),\nu)$ bilden.

\begin{center}
\textbf{§ 5. \quad Zurückführung des Moduls $(abc)$ auf die Moduln $\A$ und $\nu$.}
\end{center}

Der nur durch einen Klammerfaktor gegebene Modul $(abc)$ ist in Übereinstimmung mit der Definition (§ 3, a) auf Formen des Systems zurückzuführen. In anderen Worten: die Formenreihe $(abc)$ ist eindeutig zu normieren (vgl. § 1).

\[
a_s^2a_xu_x=\frac13 a_s^3u_x=\frac13 S\cdot u_x. \tag{2}
\]
\[
\left.
\begin{aligned}
(a\A u)^2&=a_s-\frac{S}{6}u_x^2,\\
(a\A u)^3&=0.
\end{aligned}
\right\} \tag{3}
\]
\[
\left.
\begin{aligned}
\theta(s)=(ss,x)^2&=2a_t+\frac13 S\cdot\theta-\frac13 Tu_x^2,\\
\theta(t)=(tt,x)^2&=-\frac13 S\cdot a_t+\frac23 T\cdot a_s+\frac1{18}S^2\cdot\theta-\frac1{18}u_x^2\cdot S\cdot T.
\end{aligned}
\right\} \tag{4}
\]

Reihe der Moduln: $(abc)$; $\A$; $s$; $t$ (das System des Moduls $t$ besteht aus der einzigen Form $t$).\footnote{Gordan--Kerschensteiner, a.\,a.\,O. S. 134.}

Es wird sich zeigen, daß die Formen:
\[
\A=(abc)^2a_x^2b_x^2c_x^2,\qquad \nu=(abu)^4
\]
als Modul genügen, d.\,h. daß die Formen
\[
\begin{gathered}
\A=(abc)^2a_x^2b_x^2c_x^2,\qquad
a_\nu=(abc)(bcu)^3a_x^3,\\
a_\nu^2=(abc)^2(bcu)^2u_x^2,\qquad
i=a_\nu^4=(abc)^4
\end{gathered}
\]
die Formenreihe $(abc)$ definieren. In der Formenreihe $a_\nu$ fehlt die Form $a_\nu^3$ nach der Identität:
\[
a_\nu^3=(abc)^3a_x(bcu)=\frac13 u_x\cdot(abc)^4=\frac13 i\cdot u_x.
\]

Zur Zurückführung eines Moduls auf einen höheren haben wir nach der Definition die allgemeinsten Faltungen oder auch alle speziellen in Betracht kommenden Faltungen mit dem Modul auf Faltungen mit dem höheren Modul zurückzuführen.

Die allgemeinsten Faltungen entstehen in unserm Fall aus den Ausdrücken:
\[
(abc)a_x^3b_y^3c_z^3,\qquad
(abc)a_x^2b_y^2c_z^2a_zb_xc_x
\]
(aus den kubischen Formen übernommen),
\[
(abc)a_xb_yc_za_x^2b_x^2c_x^2
\]
(aus den quadratischen Formen übernommen) und durch Polarisation dieser Ausdrücke.

Zur Berechnung dieser Ausdrücke durch die Polaren von $\A$ und der Formenreihe $a_\nu$, bzw. deren Anfangsglieder, schlagen wir, wie in § 2, den indirekten Weg ein. Wir entwickeln die Polarenglieder
\[
(abc)^2a_x^2b_y^2c_z^2,\qquad
a_\nu(\nu xy)(\nu yz)(\nu zx)a_xa_ya_z
=(abc)(bc\widehat{x}u)(bc\widehat{y}z)(bc\widehat{z}x)a_xa_ya_z\quad\text{usw.}
\]
als lineare Aggregate von Ausdrücken mit dem symbolischen Faktor $(abc)$ und erhalten durch Umkehrung des Gleichungssystems die gesuchten Relationen, sowie eine Relation zwischen den nach verschiedenen Kombinationen der Variablen genommenen Polaren. Zur Auswertung dient der Identitätssatz und der Produktsatz für Determinanten.

Das Gleichungssystem lautet für $(abc)a_x^3b_y^3c_z^3$:
\begingroup\small
\[
\begin{aligned}
1)\quad
\sum_3 a_\nu(\nu yz)^3a_x^3
&=6(abc)a_x^3b_y^3c_z^3
-6\sum_3(abc)a_x^3b_y^2b_zc_z^2c_y^{*},\\[0.4em]
2)\quad
3(abc)^2a_x^2b_y^2c_z^2\cdot(xyz)
-\frac12\sum_3 a_\nu^2(\nu yz)^2\nu_x^2(xyz)
&=2\sum_3(abc)a_x^3b_y^2b_zc_z^2c_y\\
&\quad-4\sum_3(abc)a_xa_ya_zb_y^2b_zc_z^2c_x,\\[0.4em]
3)\quad
a_\nu(\nu xy)(\nu yz)(\nu zx)a_xa_ya_z
&=2\sum_3(abc)a_xa_ya_zb_y^2b_zc_z^2c_x,\\[0.4em]
4)\quad
\sum_3 a_\nu(\nu yz)^3u_x^3
-\frac32\sum_3 a_\nu^2(\nu yz)^2\nu_x^2(xyz)
+\frac16 i\cdot(xyz)^3
&=6\sum_3(abc)a_xa_ya_zb_y^2b_zc_z^2c_x.
\end{aligned}
\]
\endgroup

Daraus folgt für
\[
(abc)a_x^3b_y^3c_z^3,
\]
und durch analoge Rechnung für
\[
(abc)a_x^2b_y^2c_z^2a_zb_xc_x,\qquad
(abc)a_xb_yc_za_z^2b_x^2c_x^2:
\]
\[
\begin{aligned}
(abc)a_x^3b_y^3c_z^3
&=\frac32(abc)^2a_x^2b_y^2c_z^2(xyz)
+\frac32 a_\nu(\nu xy)(\nu yz)(\nu zx)a_xa_ya_z\\
&\quad-\frac1{12}i\cdot(xyz)^3,\\[0.4em]
(\text{IV.})\quad
(abc)a_x^2b_y^2c_z^2a_zb_xc_x
&=\frac23(abc)^2a_x^2b_y^2c_z^2c_x\cdot(xyz)\\
&\quad+\frac13a_\nu(\nu xy)(\nu yz)(\nu zx)a_x^3
-\frac16 a_\nu^2(\nu xy)(\nu zx)a_x^2\cdot(xyz),\\[0.4em]
(abc)a_xb_yc_za_z^2b_x^2c_x^2
&=\frac16(abc)^2a_x^2b_z^2c_z^2\cdot(xyz).
\end{aligned}
\]

Wir haben somit ein relativ vollständiges System mod $(\A,\nu)$ erhalten, bestehend aus den vier Formen: $f,\theta,K,j$. Es folgt daraus: Alle nicht aus dem binären Gebiet übernommenen Überschiebungen von $f$ über $\theta$, ebenso wie die im Binären auf $(ab)^4$ reduzible Überschiebung $(a\theta u)^2$ lassen sich ausdrücken durch Symbole $\A$ und $\nu$.\footnote{$\sum_3$ bezieht sich auf die durch zyklische Vertauschung der Variablen aus dem Anfangsglied entstehende Summe über drei Glieder. Die Gleichungen gelten auch einzeln für jedes Glied der Summe.}

Wir erhalten die für das Folgende grundlegenden Reduktionsformeln, durch Spezialisierung von Entwicklung IV, oder kürzer durch direkte Rechnung (Vertauschen der Symbole), für $a_\vartheta(a\theta u)^2$, $a_\vartheta^2((a\theta u)^2$, $(a\theta u)^2$ nach folgendem Ansatz:
\[
\begin{aligned}
a_\vartheta(a\theta u)^2
&=(abc)(bcu)(abu)^2-\frac13(abc)(bcu)\{(bcu)a_x-u_x(abc)\}^2,\\
a_\vartheta^2(a\theta u)^2
&=(abc)^2(abu)^2-\frac13(abc)^2\{(bcu)a_x-u_x(abc)\}^2,
\end{aligned}
\]
für $((a\theta u)^2)^{*}$:
\[
\begin{aligned}
(abu)a_y^3b_x^3&=\frac32u_\vartheta(\vartheta yx)\theta_y^2+\frac14(\nu yx)^3,\\
((abu)(acu))b_x^3&=\frac32u_\vartheta(\vartheta\widehat{a}ux)(\theta au)^2+\frac14(\nu\widehat{a}ux)^3.
\end{aligned}
\]
\[
\left.
\begin{aligned}
(a\theta u)^2&=\frac16\nu\cdot f-\frac23u_x\cdot a_\nu
+\frac23u_x^2\cdot a_\nu^2-\frac{i}{18}u_x^4,\\
a_\vartheta&=\frac13u_x\cdot\A,\\
a_\vartheta(a\theta u)&=0,\\
a_\vartheta^2&=\A,\\
(a\theta u)^3&=0,\\
a_\vartheta(a\theta u)^2&=-\frac56a_\nu+\frac76u_x\cdot a_\nu^2-\frac{i}{9}u_x^3,\\
a_\vartheta^2(a\theta u)&=0,\\
(a\theta u)^4&=j,\\
a_\vartheta(a\theta u)^3&=0,\\
a_\vartheta^2(a\theta u)^2&=\frac23a_\nu^2-\frac{i}{9}u_x^2.
\end{aligned}
\right\} \tag{1}
\]

Wir nehmen dazu die ebenfalls aus der Reduktion des Moduls $(abc)$ entstandene Formel:
\[
a_\nu^3a_xu_\nu=\frac13 i\cdot u_x. \tag{2}
\]

\emph{Aus den Formeln (1.) und (2.) und den für die Grundform $\nu$ analog gebildeten Formeln leiten sich alle späteren Reduktionsformeln durch Polarisation ab}, ausgenommen die Relationen für die zerfallenden Formen. Aus den Formeln (1.) sehen wir:

Es führt die Formenreihe:
\[
\begin{array}{rcl}
a_\vartheta(a\theta u)^2 &\text{auf}& \text{Symbole }\nu\text{ allein},\\[0.2em]
a_\vartheta &\text{auf}& \text{Symbole }\A\text{ und }\nu,\\[0.2em]
(a\theta u)^2 &\text{auf}& \text{Symbole }j,\A\text{ und }\nu,\\[0.2em]
(a\theta u)^2\text{ bei Faltung I} &\text{auf}& \text{Symbole }j\text{ und }\nu,\\[0.2em]
(a\theta u)^2\text{ bei Faltung II} &\text{auf}& \text{Symbole }\A\text{ und }\nu.
\end{array}
\]

Wir können also ersetzen:
\begin{enumerate}
\item mod $(\A,\nu)$: Faltungen mit $j$ durch entsprechende mit $(a\theta u)^2$ oder $(a\theta u)^3$;
\item mod $(\nu)$: Faltungen mit $\A$ durch entsprechende mit $a_\vartheta$ oder $a_\vartheta(a\theta u)$ oder auch durch spezielle mit $(a\theta u)^2$ (die nur zu einmaliger Faltung I führen).
\end{enumerate}
Es wird insbesondere:
\[
\begin{aligned}
(a\theta u)^2\theta_y^2\theta_x^2&=\frac13(jyx)^2\pmod{\nu},&
(a\theta u)^2u_y\theta_y&=-\frac16(jyx)^2\pmod{\nu},\\
(a\theta u)^2\nu_y^2&=\frac13(\A u\nu)^2\pmod{\nu},&
(a\theta u)(a\theta\nu)u_\vartheta\nu_\vartheta&=-\frac16(\A u\nu)^2\pmod{\nu},\\
(a\theta u)^2u_\nu^2\vartheta_\nu^2&=\frac13(\A u\nu)^2\A_\nu\pmod{\nu}.
\end{aligned}
\]

\begin{center}
\textbf{§ 6. \quad Zurückführung des Moduls $(\A,\nu)$ auf den Modul $(\nu)$.}
\end{center}

Die Reduktionsformeln des letzten Paragraphen geben uns das Mittel, direkt das relativ vollständige System mod $\nu$ aufzustellen; wir werden sehen, daß wir dem System der übernommenen Formen nur die Formen $\A$ und $(a\A u)=N$ beizufügen haben (analog wie bei der kubischen Form und überhaupt allgemein gültig).

Wir betrachten zur Reduktion des Moduls $\A$ alle speziellen in Betracht kommenden Faltungen, d.\,h. die Faltungen des relativ vollständigen Systems mod $(\A,\nu)$ mit dem System von $\A$, und gehen aus von den in den Koeffizienten niedrigsten Formen, den Faltungen von $f$ mit $\A$.

Zur Reduktion der Formen $(a\A u)^2$, $(a\A u)^3$, $(a\A u)^4$ ersetzen wir nach § 5 die Symbole $\A$ durch niedrigere Formen der Formenreihe, d.\,h. wir entwickeln die Ausdrücke
\[
a_\vartheta b_\vartheta(b\theta u)^2,\qquad
a_\vartheta(a\theta u)b_\vartheta(b\theta u)^2,\qquad
a_\vartheta(a\theta u)b_\vartheta(b\theta u)^3
\]
1) nach Polaren der Formenreihe $a_\vartheta$ (Symbole $\A$ und $\nu$),

2) nach Polaren der Formenreihe $b_\vartheta(b\theta u)^2$ (Symbole $\nu$)

und erhalten die Reduktionsformeln (Rechnung siehe unten):
\begin{align}
\text{a)}\quad (a\A u)^2
&=\frac12(\vartheta\nu x)^2-\frac25 f\cdot a_\nu^2
-\frac45u_x\cdot\theta_\nu(\vartheta\nu x)^2 \notag\\
&\quad+u_x^2\left\{\frac16 i\cdot f-\frac1{40}S\right\},\notag\\
\text{b)}\quad (a\A u)^3
&=-\frac3{10}\theta_\nu(\vartheta\nu x),\tag{3}\\
\text{c)}\quad (a\A u)^4
&=\frac{21}{10}\theta_\nu^2-\frac3{10}\Pi
-\frac65u_x\cdot\theta_\nu^3+\frac1{10}u_x^2\cdot\theta .\notag
\end{align}

(Zu den gleichen Resultaten führt auch die doppelte Reduktion der Ausdrücke $a_\vartheta^2(b\theta u)^2$, $a_\vartheta^2(b\theta u)^3$, $a_\vartheta^2((a\theta u)^2(b\theta u)^2$ und $a_\vartheta^2((a\theta u)(b\theta u))^3$ nach Elimination von $a_\vartheta^2$.)

Aus den Formeln (3.) folgt, daß $(a\A u)^2$ Reduzent ist in bezug auf den Modul $\nu$. Daraus ergibt sich:

1) Die Reduktion des Systems von $\A$: Die Formenreihe $(\A\A' u)^2=a_\vartheta^2((a\A u)^2$ hat den Reduzenten $(a\A u)^2$ zum Faktor.

2) Die Reduktion des durch Faltung mit dem Modul $\A$ entstandenen Systems:

Die Formenreihe $(\theta\A u)^2$ hat den Reduzenten $(a\A u)^2$ zum Faktor.

Für die Formen $(\theta\A u)$ und $\A_\vartheta$ ergibt sich:
\[
(a\A u)^2(abu)=(\theta\A u)-u_x\cdot\A_\vartheta(\theta\A u),
\]
\[
(a\A u)(a\A b)=-\frac12\A_\vartheta+\frac12u_x\cdot\A_\vartheta^2.
\]

Aus dem Reduzenten $(\theta\A u)$ folgt aber die Reduktion des durch Faltung mit dem Modul $\A$ entstehenden Systems.
\end{document}
